% !TEX program = xelatex
\documentclass{英语试卷}
\begin{document}
\examtitle{温州市普通高中 2025 届高三第三次适应性考试}{英语}
\par 温州市普通高中 2025 届高三第三次适应性考试\par
\par 英语试题\par
\begin{notice}
\par 1. 答卷前，考生务必将自己的姓名、准考证号填写在答题卡上。\par
\par 2. 回答选择题时，选出每小题答案后，用铅笔把答题卡上对应题目的答案标号涂黑。如需改动，用橡皮擦干净后，再选涂其他答案标号。回答非选择题时，将答案写在答题卡上，写在本试卷上无效。\par
\par 3. 考试结束后，将本试卷和答题卡一并交回。\par
\end{notice}
\parttitle{第一部分 听力（共两节，满分30分）}
\par 做题时，先将答案标在试卷上，录音内容结束后，你将有两分钟的时间将试卷上的答案转涂到答题卡上。\par
\sectitle{第一节（共5小题：每小题1.5分，满分7.5分）微}
\par 听下面 5 段对话。每段对话后有一个小题，从题中所给的 A、B、C 三个选项中选出最佳选项。听完每段对话后，你都有 10 秒钟的时间来回答有关小题和阅读下一小题。每段对话仅读一遍。\par
\question{1. What is Megan doing?}
\begin{choices}(3)
  \item Doing her homework.
  \item Preparing for a trip
  \item Changing her clothes.
\end{choices}
\question{2. Why does the man suggest starting a book-sharing club?}
\begin{choices}(3)
  \item It reduces costs.
  \item It saves time.
  \item It protects nature
\end{choices}
\question{3. What does the woman do?}
\begin{choices}(3)
  \item She's a hotel manager
  \item She's a swimming coach.
  \item She's a gym receptionist.
\end{choices}
\question{4. What has caused the power failure?}
\begin{choices}(3)
  \item A heavy storm.
  \item The unpaid bill.
  \item High powered devices.
\end{choices}
\question{5. Where are the speakers?}
\begin{choices}(3)
  \item At a science fair.
  \item In a geography class.
  \item At a birthday party.
\end{choices}
\sectitle{第二节 (共 15 小题: 每小题 1.5 分, 满分 22.5 分)}
\par 听下面 5 段对话或独白。每段对话或独白后有几个小题，从题中所给的 A、B、C 三个选项中选出最佳选项。听每段对话或独白前，你将有时间阅读各个小题，每小题 5 秒钟；听完后，各小题将给出 5 秒钟的作答时间。每段对话或独白读两遍。\par
\par 听第 6 段材料，回答第 6、7 题。\par
\question{6. Why does the man fly to Seattle?}
\begin{choices}(3)
  \item To attend a meeting.
  \item To visit a friend.
  \item To go sightseeing.
\end{choices}
\question{7. What will the man do next?}
\begin{choices}(3)
  \item Cancel the trip.
  \item Book a hotel.
  \item Take another flight.
\end{choices}
\par 听第 7 段材料，回答第 8-10 题。\par
\question{8. What is the possible relationship between the speakers?}
\begin{choices}(3)
  \item Classmates.
  \item Workmates.
  \item Neighbors.
\end{choices}
\question{9. How will Sarah go to Oxfordshire?}
\begin{choices}(3)
  \item By train.
  \item By plane
  \item By car.
\end{choices}
\question{10. What should Sarah do before setting out?}
\begin{choices}(3)
  \item Clean the office.
  \item Buy new plants.
  \item Buy a pet carrier.
\end{choices}
\par 听第 8 段材料，回答第 11-13 题。\par
\question{11. What does the man think of Uber Eats?}
\begin{choices}(3)
  \item Convenient.
  \item Entertaining.
  \item Complex
\end{choices}
\question{12. What does Uber Eats provide?}
\begin{choices}(3)
  \item Healthy eating tips
  \item Free delivery service.
  \item Tailored recommendations.
\end{choices}
\question{13. What does the man suggest the woman do?}
\begin{choices}(3)
  \item Order less food.
  \item Delete the app.
  \item Take some exercise.
\end{choices}
\par 听第 9 段材料，回答第 14-17 题。\par
\question{14. Who organized the competition?}
\begin{choices}(3)
  \item A film studio.
  \item A TV network.
  \item A travel agency.
\end{choices}
\question{15. What attracted the man to enter the competition?}
\begin{choices}(3)
  \item A job offer.
  \item The prize money.
  \item The professional guidance.
\end{choices}
\question{16. What does the man say about the trip?}
\begin{choices}(3)
  \item It is totally free.
  \item It is a guided tour.
  \item It lasts for three months.
\end{choices}
\question{17. What will the man's documentary focus on?}
\begin{choices}(3)
  \item Historical stories.
  \item Tourism development.
  \item Environmental practices.
\end{choices}
\par 听第 10 段材料，回答第 18-20 题。\par
\question{18. How long does the festival last?}
\begin{choices}(3)
  \item 6 hours.
  \item 8 hours
  \item 10 hours.
\end{choices}
\question{19. What do we know about the activities of the festival?}
\begin{choices}(3)
  \item They are wide-ranging.
  \item They are culture-focused.
  \item They are technology-based.
\end{choices}
\question{20. Why does the speaker give the talk?}
\begin{choices}(3)
  \item To discuss a plan.
  \item To do a promotion.
  \item To ask for donations.
\end{choices}
\parttitle{第二部分 阅读理解（共两节，满分 50 分）}
\sectitle{第一节 (共 15 小题: 每小题 2.5 分, 满分 37.5 分)}
\par 阅读下列短文，从每题所给的 A、B、C 和 D 四个选项中，选出最佳答案。\par
\passagelabel{A}
\par Power of Public Speaking has partnered with ASI Innovation Labs to integrate artificial intelligence into the Pops Classroom platform for improving public speaking skills. The first project is the Pops Advisor which provides an automatic evaluation of students' video presentations with visual feedback.\par
\par The Advisor offers on-demand comments to students practicing public speaking from anywhere. This flexibility benefits instructors, who can be as hands-on or hands-off during practice. While students can operate autonomously with the tool, instructors can also choose to keep track of students' progress. Criteria\par
\par Using machine learning models, the Pops Advisor analyzes the following criteria:\par
\par\vspace{0.4em}
\begin{center}\begin{tabular}{{|l|l|l|l|}}
\hline
Delivery & Introduction & Body & Conclusion \\
Vocal Style & Credibility & Support & Signaling Ending \\
Vocal Quality & Attention & Language & Summary \\
Vocal Filler & Central Idea & Connective & Memorable Ending \\
Eye Contact & Preview & Credibility & Credibility \\
Body Language & Connective & Organization &  \\
\hline
\end{tabular}\end{center}
\vspace{0.4em}
\par Feedback\par
\par The presenter receives feedback on the above four categories. And this feedback is given in the following formats:\par
\par - Recommendations\par
\par To provide presenters with valuable feedback for the next practice, recommendations focus on specific aspects needing improvement. For instance, the recommendation can be "Your message is clear, but be sure you are grabbing your audience's attention in your introduction to ensure they are listening."\par
\par - Speech Metrics (指标)\par
\par Useful speech metrics serve as quantitative descriptors on several aspects of the presentation including number of words, total presentation time, speaking rate and vocal fillers percentage.\par
\par Pops Advisor can be used in any classroom setting where students are presenting, particularly in the Basic Public Speaking Course. When you are considering what can be achieved by your students toward skill development in a single school semester, Advisor increases the overall development.\par
\question{21. What is the Pops Advisor?}
\begin{choices}(2)
  \item A voice assistant.
  \item A video player.
  \item A time tracker.
  \item An evaluation tool.
\end{choices}
\question{22. Which of the following categories provides feedback on facial expressions?}
\begin{choices}(2)
  \item Delivery.
  \item Introduction.
  \item Body.
  \item Conclusion.
\end{choices}
\question{23. Who is the text intended for?}
\begin{choices}(2)
  \item Lab technicians.
  \item Public speech instructors
  \item Foreign language learners.
  \item Artificial Intelligence engineers
\end{choices}
\passagelabel{B}
\par The boat ride from Helsinki's city center to the island of Pihlajasaari takes only 10 minutes and drops off visitors at beaches for sunbathing. But I had a different mission: to speak to a tree.\par
\par This was meant to be a healing exercise, one championed by the Finnish biologist Adela Pajunen. Finns, she'd told me, sometimes share their worries aloud to trees or birds. On shore, I spotted a short black alder tree and began to tell the tree my troubles. The leaves rustled (沙沙作响) in response, a sign I interpreted as sympathy\par
\par I had come to Finland to see whether I could bring happiness back to America with me. Finland has topped the World Happiness Report for the past eight years, largely due to free education and universal health care. And the Finns also find happiness in more attainable ways, such as their close relationship with nature and visiting the sauna (桑拿) daily. These activities contain similar goals: Stay present and seek a minimalist lifestyle that relies on the earth.\par
\par I followed their lifestyle—visiting the sauna and taking a relaxing forest walk, and happiness came in waves over the next few days. I was optimistic that I could recreate the happiest moments of this trip back home in New York. How hard could it be, really?\par
\par As it turned out, happiness was a luxury in America, though I tried hard to live exactly like the Finns. Then one afternoon, I had a breakthrough as I rushed through Central Park and noticed a crowd of people staring at the ground. A cardinal! I stopped to admire the bird's feather before it flew away. This was a win for staying present. I'd been so stuck trying to find the time and money for happiness that I'd missed the point: Quiet the mind, and find pleasure in small acts and observations—like a cardinal taking flight, or talking to trees—right then and there.\par
\par As I left the park, I carried out a small but comforting ritual (仪式) I'd developed. I thanked the trees for listening.\par
\question{24. Why did the author speak to a tree?}
\begin{choices}(2)
  \item To heal the tree.
  \item To seek happiness.
  \item To follow a tradition.
  \item To appreciate nature.
\end{choices}
\question{25. What do we know about the Finnish lifestyle?}
\begin{choices}(2)
  \item It requires daily sauna visits.
  \item It values close relationships
  \item It is easily adopted elsewhere.
  \item It focuses on simple pleasures
\end{choices}
\question{26. Why couldn't the author recreate happy moments back in America?}
\begin{choices}(2)
  \item She was busy with work
  \item It was noisy in New York
  \item She forgot to stay present
  \item It was expensive to live there.
\end{choices}
\question{27. Which of the following best describes the author's trip to Finland?}
\begin{choices}(2)
  \item Transformative.
  \item Challenging
  \item Adventurous.
  \item Unexpected.
\end{choices}
\passagelabel{C}
\par One of the top sources of added sugar in children's diets in America is in their breakfast cereal (麦片). A study published in the American Journal of Preventive Medicine shows that advertising drives sales of high-sugar cereals when it's aimed directly at kids under 12—but not when it targets adults.\par
\par In the study, researchers looked at all cereals purchased by 77,000 U. S. households over a nine-year period, between 2008 and 2017. They also looked at Nielsen ratings data, which closely monitored all the ads that people in a household saw—both children and adults. What they found was a strong relationship between how much advertising was targeted to kids and how much sugary cereal that households with children bought. By contrast, there was no link to increased purchases when ads targeted adults. And these extensive data also showed that behaviors that were learned in childhood could track into adulthood, which could lead to poor health outcomes over a lifetime.\par
\par The study is the first to directly compare the influence of food advertising exposure by children with that of adults. The findings offer novel evidence of how food marketing turns children into the "ultimate weapon" in influencing family spending because they consistently ask their parents to purchase—a concept known as "pester power".\par
\par The food industry launched the Children's Food and Beverage Advertising Initiative, a voluntary effort to police itself. The 21 participating food companies made a commitment to cut back on marketing unhealthy foods to children under 12—later revised to under 13. A new study conducted in 2024 showed that children's exposure to cereal ads on TV programming aimed at kids has dropped dramatically.\par
\par But Lindsey Smith Taillie, a food policy researcher at the University of North Carolina at Chapel Hill, says these voluntary efforts aren't making a difference. Advertisements, like kids' eyeballs, are moving online, making it even harder to know what marketing children are being exposed to, Taillie notes.\par
\question{28. What did the study find about child-targeted food ads?}
\begin{choices}(2)
  \item They had no connection with adults' purchases.
  \item They were less effective than adult-targeted ads.
  \item They might affect children's health in the long run.
  \item They could reduce children's interest in sugary food.
\end{choices}
\question{29. What does the underlined phrase "pester power" in paragraph 3 mean?}
\begin{choices}(2)
  \item Parents' unconditional love.
  \item Children's uncontrolled spending
  \item Parents' growing consumption
  \item Children's repeated purchase requests
\end{choices}
\question{30. What is paragraph 4 mainly about?}
\begin{choices}(3)
  \item An expectation for further study.
  \item A limitation of the research method. 3. The follow-up influence of the study
  \item The potential application of the findings.
\end{choices}
\question{31. What is Tallie's attitude toward the voluntary efforts?}
\begin{choices}(2)
  \item Favorable.
  \item Tolerant.
  \item Disapproving.
  \item Reserved
\end{choices}
\passagelabel{D}
\par Surrounded by towering redwoods on one side and the Pacific Ocean on the other, the coastal city of Santa Cruz, some 70 miles south of San Francisco, is a great place for a young person to experience nature. That's where Xinning Zhang, assistant professor of geosciences and the Princeton Environmental Institute, spent her childhood. "Growing up there, you just can't help thinking and caring about the environment," she said. One of the environmental questions Zhang is exploring is why methane, a significant greenhouse gas, is increasing in the atmosphere.\par
\par Most of us are aware that cows release methane, but wetlands are also a major source of the gas. Wetlands host a community of microbes (微生物), including Archaea, which generate methane as a byproduct when they consume and digest components of organic matter. These microbes grow in oxygen-free environments like wet soils, which typically occur in the wetlands' deeper levels. However, researchers have been surprised to find that quite a lot of methane comes from the oxygen-rich peat (泥炭) and soil near the surface.\par
\par To explore this phenomenon, Zhang and her team analyzed peat samples collected from a wetland in the Northeast U. S. One set of samples received oxygen treatment followed by incubation (培养) under oxygen-free conditions. The other samples remained in an oxygen-free environment for the entire study period. Contrary to expectations, the oxygen-treated peats produced a higher amount of methane than the peats maintained under continuously oxygen-free conditions.\par
\par To find out why, the researchers analyzed all the microbes found in the peat samples. They found that oxygen stimulated some microbes to break down harmful chemicals, which would otherwise threaten the entire microbial community, including methane-producing Archaea. The result is the creation of far more methane than expected.\par
\par Zhang's work suggests that oxygen variability is an important control on wetland methane production. Her team is also exploring further how oxygen variability for different lengths of time, and how different soil chemistries influence microbial methane production.\par
\question{32. What motivated Zhang to study environmental issues?}
\begin{choices}(2)
  \item Her childhood experiences
  \item The adventure in nature.
  \item Her educational background.
  \item The pressure from fellows.
\end{choices}
\question{33. Why did Zhang analyze microbes in the wetland soil?}
\begin{choices}(2)
  \item To examine the condition of surface soil.
  \item To explore the impact of methane on microbes
  \item To study the role of oxygen in methane production.
  \item To calculate the amount of methane released by cows.
\end{choices}
\question{34. How might the findings help reduce methane from wetlands?}
\begin{choices}(2)
  \item By changing soil structures.
  \item By introducing new microbes.
  \item By blocking harmful chemicals
  \item By adjusting soil oxygen levels.
\end{choices}
\question{35. Which of the following can be the best title for the text?}
\begin{choices}(2)
  \item How Does Methane Affect Climate Change?
  \item What Drives Methane Production in Wetlands?
  \item What Is the Link Between Oxygen and Methane?
  \item How Do Wetlands Contribute to Microbial Diversity?
\end{choices}
\sectitle{第二节（共5小题；每小题2.5分，满分12.5分）}
\par 阅读下面短文，从短文后的选项中选出能填入空白处的最佳选项。选项中有两项为多余选项。\par
\par Parenting is a constant balancing act. We want to teach our children important life lessons, while also maintaining harmony at home and avoiding unnecessary conflicts. \blanknum{36}\par
\par Recently, my 13-year-old daughter asked me to buy her a \$28 phone case. Her old one wasn't broken or worn out; she was just bored of it. \blanknum{37} However, saying "no" in moments like this isn't easy. The request felt small, but I knew there was a bigger lesson to teach. It's about teaching our children to think critically about their spending. At the same time, we don't want to come across as overly strict or dismissive of their feelings.\par
\par In this case, I realized my daughter's request wasn't just about the phone case; it was about her desire for something new, trendy, and exciting. \blanknum{38} After all, to a 13-year-old, a new phone case might feel like a big deal, a way to express themselves or fit in with their friends. So how do we find the middle ground?\par
\par Instead of shutting down the conversation, I used it as an opportunity to talk about wants and needs. I explained that we must prioritize what's essential over impulse (冲动) purchases. \blanknum{39} I also acknowledged her feelings. "I get it," I told her. Sometimes we just want something new because it feels exciting. This softened the conversation and made her more willing to listen. Instead of a flat-out "no", I suggested a compromise (妥协). I offered her three options: earn it, wait it out or explore alternatives.\par
\par \blanknum{40} It was about guiding my daughter toward better decision-making without damaging our relationship.\par
\begin{choices}(1)
  \item But I didn't stop there
  \item For me, it was a firm no
  \item That conversation wasn't just about a phone case.
  \item I promoted the habit of saving for things she wanted.
  \item It felt important to her, even if it seemed unnecessary to me.
  \item This balancing act can be tricky, especially when parenting teens.
  \item Sometimes, the most valuable lessons aren't about the money at all.
\end{choices}
\parttitle{第三部分 语言知识运用（共两节，满分 30 分）}
\sectitle{第一节 (共 15 小题; 每小题 1 分, 满分 15 分)}
\par 阅读下面短文，从短文后各题所给的 A、B、C、D 四个选项中，选出可以填入空白处的最佳选项.\par
\par I love everything about our restaurant La Salle Rouge. In the kitchen, the crash of forks and knives hitting the heavy plates almost \blanknum{41} Mom's voice calling me. "Lainey? Lai-ney!" \blanknum{42}, I look up to see my mother standing across the kitchen. "Homework?" Mom mouths the word, pointing to the side door. I roll my eyes—I \blanknum{43} the thought of leaving the kitchen to \blanknum{44} with my math homework in Mom's office.\par
\par The restaurant has always been the \blanknum{45} of our lives. At fourteen, I started being a vegetarian (素食者), but Mom still found things to \blanknum{46} me and taught me to be creative with vegetables and tofu. When I started high school thirty pounds heavier than everyone in my class, Mom and I designed a \blanknum{47} menu for La Salle Rouge, and Mom let me \blanknum{48} tasty, low-calorie desserts, which became popular later. I \blanknum{49} early on that this was the work I want to do.\par
\par It hardly seems \blanknum{50} that I have to walk away from all of that just to do maths, but Mom says it's smarter to have a backup \blanknum{51} and she's made me apply to colleges and check out business majors \blanknum{52} I want to do anything else with my life. But I already know what I'm going to do when I graduate.\par
\par First, I'll buy myself a set of knives, which is a must for the best chefs. The second thing is to get \blanknum{53}—people would know my name and try my \blanknum{54} Now, all I have to do is finish homework and get back into the \blanknum{55}\par
\cz{41}{attracts}{drowns}{raises}{reduces}
\cz{42}{Sighing}{Choking}{Shaking}{Nodding}
\cz{43}{accept}{prefer}{change}{hate}
\cz{44}{wrestle}{argue}{experiment}{play}
\cz{45}{truth}{standard}{duty}{center}
\cz{46}{help}{feed}{save}{control}
\cz{47}{traditional}{romantic}{light}{diverse}
\cz{48}{cut down on}{come up with}{run out of}{get rid of}
\cz{49}{promised}{suggested}{decided}{predicted}
\cz{50}{fair}{possible}{strange}{enough}
\cz{51}{tool}{system}{plan}{design}
\cz{52}{as though}{now that}{in case}{even if}
\cz{53}{paid}{married}{fired}{discovered}
\cz{54}{books}{recipes}{routines}{products}
\cz{55}{kitchen}{school}{house}{office}
\sectitle{第二节 (共 10 小题; 每小题 1.5 分, 满分 15 分)}
\par 阅读下面短文，在空白处填入 1 个适当的单词或括号内单词的正确形式。\par
\par Blind boxes have taken the world by storm, catching the imaginations of collectors and new buyers alike. These mystery boxes, containing \blanknum{} \blanknum{56}\blanknum{} (hide) items revealed only after purchase, offer a mix of surprise and exclusivity (独家性). From toys to high-end collectibles, \blanknum{} \blanknum{57}\blanknum{} blind box trend attracts anyone seeking excitement or wanting to add something unique to their collection.\par
\par The appeal lies in the excitement of not knowing \blanknum{58} is inside, combined with the desire to complete a set or find rare items. This concept of discovery keeps collectors \blanknum{59} (come) back for more and has helped transform the blind box industry into a multibillion-dollar market.\par
\par For those new \blanknum{60} the blind box world, starting a collection can be a rewarding experience. Begin by selecting a brand or theme that matches your interests. Many collectors recommend focusing on a specific series, which makes the collection more personal \blanknum{61} helps prevent unnecessary purchases. Engaging with the community can also enhance your collecting journey. There are countless social media groups where collectors share tips and discuss new \blanknum{62} (release). Connecting with other collectors \blanknum{63} (allow) you to trade the same items, find rare pieces, and share in the excitement of each new \blanknum{64} (add) to your collection.\par
\par Whether you're an experienced collector or a newcomer, the world of blind boxes offers \blanknum{65} (end) possibilities, turning every unboxing into an adventure.\par
\parttitle{第四部分 写作（共两节，满分 40 分）}
\sectitle{第一节 (满分 15 分)}
\par 假定你是李华，你校计划举办“时光机音乐节 (Time Machine Music Festival)”，你将与英国交换生 Tom 共演一个节目。请你写一封邮件与他商讨，内容包括：\par
\par (1) 节目设想;\par
\par (2) 设计理由;\par
\par (3) 征求意见。\par
\par 注意:\par
\par (1) 写作词数应为 80 左右;\par
\par (2) 可适当增加细节，以使行文连贯。\par
\par\vspace{0.4em}
\begin{center}\begin{tabular}{{|l|}}
\hline
Dear Tom, \\
Yours, Li Hua \\
\hline
\end{tabular}\end{center}
\vspace{0.4em}
\sectitle{第二节 (满分 25 分)}
\par 阅读下面材料，根据其内容和所给段落开头语续写两段，使之构成一篇完整的短文。\par
\par It was June, the last day of school. My class of first-graders was feverishly anticipating their summer vacation. To add a joyful flavor, I had some balloons brought in, but for me, this holiday was bittersweet. I'd just moved to this city and couldn't afford to fly back to accompany Dad. He was desperately missing me, for it was his first summer without me\par
\par When switching off the light on my way out, I noticed the balloons near the window. Unsure what to do with them, I grabbed the thirteen strings and stuffed them into the backseat of my car. Their rainbow of colors lifted my spirits and gratefully, I didn't feel so alone.\par
\par As the city was still foreign to me, I decided to go for a drive and soon found myself winding down a stretch of road where the trees overhead covered each other. The branches reminded me of two friends stretching out to greet each other. Before long, a retirement home appeared. It was a place with a feeling of forgottenness and I spotted a man sitting alone in the shade outside.\par
\par I felt a sudden urge to pull in. But the balloons in my backseat blocked my view—then I knew what to do with them. Opening the door and grabbing one of the strings, I waved to the man, and he smiled. His name was Daniel and he was about Dad's age. When I handed him the balloon, he grinned (咧嘴笑), "Thanks. It gives a little color to my life." While tying it to his chair, he looked happier than when I arrived. So I smiled, "You’re welcome, Daniel. Perhaps I’ll see you again sometime."\par
\par Unexpectedly, he replied, "Perhaps, but probably not. Visitors are always welcome, but rarely come." His words forced me to slide over an old chair next to his. Then we began chatting about our lives and his voice always reminded me of Dad. I'd thought it would be boring talking with a stranger. But it wasn't; it was enlightening—I had been feeling lonely, and so had Daniel, and so had the residents in the retirement home. 注意:\par
\par (1) 续写词数应为 150 左右;\par
\par (2) 请按如下格式在答题卡的相应位置作答.\par
\writingpara{Paragraph 1: Then I remembered the other twelve balloons in the backseat. \blanknum{}}
\writinglines{8}
\writingpara{Paragraph 2: Over the next six months, I visited Daniel often. \blanknum{}}
\writinglines{8}
\end{document}
